% Auto-generado por scripts/extract-lessons-to-tex.ps1
% Fuente: HTML en carpeta L3-Reconfiguracion-dogmatica-del-injusto-y-la-culpabilidad (sin modificar archivos originales).

\documentclass[11pt,a4paper]{article}
\usepackage[utf8]{inputenc}
\usepackage[T1]{fontenc}
\usepackage[spanish]{babel}
\usepackage{geometry}
\geometry{margin=2.5cm}
\usepackage{hyperref}
\hypersetup{colorlinks=true, linkcolor=blue, urlcolor=blue}

\title{\texttt{L3-Reconfiguracion-dogmatica-del-injusto-y-la-culpabilidad}}
\author{Extracci\'on desde lecciones HTML}
\date{}

\begin{document}
\maketitle

\begin{small}
\noindent Este documento consolida el texto visible de los ocho componentes (01--08). Los formularios interactivos aparecen como texto; la l\'ogica JavaScript no se reproduce aqu\'i.
\end{small}
\bigskip
\section{Conceptos}
\noindent\textit{Archivo fuente:} \texttt{01-conceptos.html}

L3 — Conceptos\par





    L3 — Elemento 1 — Conceptos\par
    El injusto penal y la culpabilidad\par

    Mapa dogmático: conducta → tipicidad → antijuridicidad → culpabilidad (esquema clásico reconfigurado en clave constitucional).\par


      La reconfiguración del injusto en el derecho penal contemporáneo integra garantías constitucionales: la tipicidad opera como\par
      reserva de ley (L1); la antijuridicidad delimita la lesión o puesta en peligro del bien jurídico; la culpabilidad funda el\par
      reproche personal y condiciona la pena y su dosimetría (L2).\par



      1. Conducta base\par

      Humana, voluntaria (en los términos que el tipo exija), orientada a resultado o creación de riesgo según el delito. Sin conducto relevante no hay examen de injusto.\par




      2. Tipicidad\par


        Ajuste de la conducta al tipo penal. Implica subsumición bajo lex certa y lex stricta. La discusión\par
        típica es el primer filtro de intervención estatal grave.\par





      3. Antijuridicidad / lesividad\par


        La conducta típica debe ser antijurídica: lesión o puesta en peligro del bien jurídico protegido, salvo causas de justificación\par
        que eliminen o reduzcan el injusto material según el caso.\par





      4. Culpabilidad\par



          DimensiónPreguntaFunción\par




          Imputabilidad¿Puede el sujeto ser sujeto de reproche?Excluyentes de responsabilidad según ordenamiento\par

          Reproche (dolo / culpa)¿Qué actitud mental exige el tipo?Personalización de la infracción\par

          Error¿Obstaculiza el conocimiento del injusto?Puede excluir o modular dolo/culpa\par








      5. Autocomprobación\par

      En el esquema clásico del injusto, la culpabilidad se examina:\par

      Antes de la tipicidad\par
      Tras la antijuridicidad\par





      Fuentes: Temario L3. Bernal, Cap. 4 (tipicidad, acto, lesividad, culpabilidad; pp. \textasciitilde{}366–418).\par


    Siguiente →\par

\section{Explicaci\'on}
\noindent\textit{Archivo fuente:} \texttt{02-explicacion.html}

L3 — Explicación\par





    L3 — Elemento 2 — Explicación\par
    Encadenamiento: de la lesividad al reproche\par


      La lesividad conecta el resultado o el riesgo con el bien jurídico tutelado. La culpabilidad traduce esa\par
      antijuridicidad en responsabilidad personal: sin culpabilidad no hay pena (principio de culpabilidad como corolario de dignidad y legalidad).\par



      1. Tipicidad y bien jurídico\par


        Cada tipo protege bienes jurídicos diferenciados. La interpretación del tipo debe preservar la función protectora sin disolver la\par
        claridad normativa (L1). La subsunción incorrecta afecta toda la cadena hasta la culpabilidad.\par





      2. Culpabilidad y pena\par


        Principio de culpabilidadSolo se sanciona culpablemente a quien es imputable y actúa con dolo o culpa según el tipo.\par

        Conexión con dosimetría (L2)La gravedad del dolo, la forma de culpa y las circunstancias morales informan la pena dentro del marco legal.\par






      3. Errores y exclusiones\par



          FiguraEfecto típico (orientativo)Nota\par




          Error de tipoPuede excluir dolo según casoRequiere análisis de conocimiento del hecho típico\par

          Error de prohibiciónPuede modular responsabilidadDistinguir invencible / vencible según sistema\par

          Causas de justificaciónExcluyen antijuridicidadNo llega a culpabilidad si no hay injusto\par





      Ajuste cada conclusión al Código Penal y jurisprudencia vigentes aplicables al caso concreto.\par




      Fuentes: Bernal, Cap. 4 (culpabilidad, tipicidad, lesividad).\par


    Siguiente →\par

\section{Exploraci\'on}
\noindent\textit{Archivo fuente:} \texttt{03-exploracion.html}

L3 — Exploración\par





    L3 — Elemento 3 — Exploración\par
    Dilema: culpabilidad “estricta” frente a políticas de prevención\par


      Algunos discursos políticos presionan por responsabilidad objetiva o inversión de la carga probatoria en ciertos delitos.\par
      La exploración contrasta esas propuestas con el principio de culpabilidad y el derecho de defensa.\par



       A) La prevención justifica atenuar exigencias de culpabilidad en delitos graves.\par
       B) La culpabilidad es límite constitucional: solo penas sobre base de imputación personal fundada.\par
       C) La tipicidad basta; la culpabilidad es formalidad accesoria.\par
      Retroalimentación\par





      Tabla argumentativa\par



          Postura preventiva extremaContrapeso garantista\par




          “El fin de disuadir autoriza presumir dolo”Inocencia y carga de la prueba; motivación racional\par

          “Basta el daño para castigar”Tipicidad, antijuridicidad y culpabilidad como filtros acumulativos\par







    Fuentes: Bernal, Cap. 4; Constitución (debido proceso, inocencia).\par

    Siguiente →\par

\section{Contexto}
\noindent\textit{Archivo fuente:} \texttt{04-contexto.html}

L3 — Contexto\par





    L3 — Elemento 4 — Contexto\par
    Constitucionalización del injusto y estándar probatorio\par


      El juzgamiento penal opera bajo debido proceso, presunción de inocencia y derecho de defensa.\par
      Cada eslabón del injusto debe reflejarse en la sentencia con estándar beyond reasonable doubt / convicción racional según el proceso aplicable.\par





          EslabónPrueba típicaError judicial frecuente\par




          TipicidadPrueba de conducta y encaje normativoConfundir indicios con subsumición completa\par

          AntijuridicidadResultado, riesgo, causas de justificaciónOmitir análisis de excluyentes\par

          CulpabilidadPrueba de dolo/culpa, errores, imputabilidadInferir dolo desde la gravedad del resultado\par







    Fuentes: Bernal, Cap. 4; Constitución arts. 28–29.\par

    Siguiente →\par

\section{Aplicaci\'on}
\noindent\textit{Archivo fuente:} \texttt{05-aplicacion.html}

L3 — Aplicación\par





    L3 — Elemento 5 — Aplicación\par
    Clínica: alegato sobre ausencia de culpabilidad / error\par


      Estructura: (1) hechos probados sobre conocimiento y voluntad; (2) estándar del tipo subjetivo; (3) narración del error o exclusión de imputabilidad;\par
      (4) petitorio de absolución o reconducción típica.\par



      Tesis modelo\par



-- La fiscalía no probó más allá de duda razonable el elemento subjetivo exigido por el tipo (indicar artículo).\par


-- Subsiste un error de tipo / prohibición vencible o invencible que excluye o modula el dolo según el caso.\par


-- La sentencia incurrió en error in iudicando al inferir dolo desde la gravedad objetiva del resultado.\par







      Verificar\par




    Fuentes: Bernal, Cap. 4; C.P. y ley procesal aplicable.\par

    Siguiente →\par

\section{Prospectiva}
\noindent\textit{Archivo fuente:} \texttt{06-prospectiva.html}

L3 — Prospectiva\par





    L3 — Elemento 6 — Prospectiva\par
    Escenarios: tipos dolosos híbridos, delitos de peligro y prueba digital\par


      La dogmática del injusto enfrenta nuevos tipos con elementos normativos densos y prueba basada en metadatos. La culpabilidad\par
      exigirá mayor precisión en imputación subjetiva cuando la conducta sea mediada por algoritmos o terceros.\par





          TendenciaRiesgo dogmáticoRespuesta técnica\par




          Tipos con “dolo eventual” discutidoConfusión con culpa graveAnálisis fino de probabilidad y aceptación del riesgo\par

          Delitos de peligro abstractoTipicidad sin lesión acreditadaControl de constitucionalidad y proporcionalidad\par

          Prueba digital masivaInferencias de dolo desde correlacionesContradicción pericial y estándar de certeza\par







    Fuentes: Bernal, Cap. 4.\par

    Siguiente →\par

\section{Ejercicios}
\noindent\textit{Archivo fuente:} \texttt{07-ejercicios.html}

L3 — Ejercicios\par





    L3 — Elemento 7 — Ejercicios\par
    Ejercicios con modelo\par


      1. Defina injusto penal en el esquema de cuatro pasos.\par

      Modelo\par





      2. ¿Qué relación hay entre antijuridicidad y culpabilidad?\par

      Modelo\par





      3. Explique por qué la gravedad del resultado no prueba automáticamente el dolo.\par

      Modelo\par





      4. Diferencie error de tipo y error de prohibición en una línea cada uno.\par

      Modelo\par





      5. ¿Qué efecto tiene una causa de justificación plena sobre la cadena del injusto?\par

      Modelo\par





      6. Vincule principio de culpabilidad con principio de dignidad humana (L1).\par

      Modelo\par




    Fuentes: Bernal, Cap. 4.\par

    Quiz →\par

\section{Quiz}
\noindent\textit{Archivo fuente:} \texttt{08-quiz.html}

L3 — Quiz\par





    L3 — Elemento 8 — Quiz\par
    Evaluación L3: injusto y culpabilidad\par

    12 preguntas con retroalimentación.\par


      1. El primer filtro dogmático suele ser:\par

         A) Culpabilidad.\par
         B) Tipicidad.\par
         C) Ejecución de la pena.\par


      2. La antijuridicidad presupone normalmente:\par

         A) Conducta atípica.\par
         B) Conducta típica.\par
         C) Solo dolo preterintencional.\par


      3. Sin culpabilidad, en el esquema clásico:\par

         A) Puede imponerse pena si hay peligro abstracto.\par
         B) No hay fundamento para pena.\par
         C) Solo cabe multa administrativa.\par


      4. Una causa de justificación excluye:\par

         A) Tipicidad.\par
         B) Antijuridicidad (injusto objetivo).\par
         C) La nacionalidad del imputado.\par


      5. Inferir dolo solo desde la gravedad del resultado es:\par

         A) Método dogmáticamente aceptable sin más.\par
         B) Cuestionable: el subjetivo requiere prueba específica.\par
         C) Obligatorio en delitos de resultado.\par


      6. El error de tipo versa sobre:\par

         A) Equivocación sobre hechos que configuran el tipo.\par
         B) Solo la pena aplicable.\par
         C) La competencia del juez.\par


      7. La tipicidad está íntimamente ligada a:\par

         A) Lex certa y reserva de ley.\par
         B) Solo la ejecución penal.\par
         C) El presupuesto administrativo.\par


      8. La imputabilidad se relaciona con:\par

         A) Capacidad de ser sujeto de reproche.\par
         B) Solo la edad del juez.\par
         C) El trámite de notificación.\par


      9. La presunción de inocencia exige:\par

         A) Que la fiscalía pruebe los elementos del delito según el estándar constitucional y legal.\par
         B) Que el acusado pruebe su inocencia siempre.\par
         C) Que no haya víctima.\par


      10. La lesividad del bien jurídico se discute en el plano de:\par

         A) Tipicidad y antijuridicidad.\par
         B) Solo la apelación de honorarios.\par
         C) Solo medidas cautelares civiles.\par


      11. El principio de culpabilidad limita:\par

         A) La pena sin conexión con reproche personal.\par
         B) Solo las multas de tránsito.\par
         C) La investigación preprocesal únicamente.\par


      12. En política criminal extrema, atenuar culpabilidad por gravedad del hecho:\par

         A) Es siempre constitucional.\par
         B) Choca con garantías de inocencia y tipicidad estricta.\par
         C) Es mandatorio.\par


      Enviar\par


      Resultado\par









    Inicio L3\par

\end{document}

